\section{Repository and How to Run the Code}

\begin{tabularx}{\textwidth}{@{}lL@{}}
  \textbf{GitHub username} & Graf-J \\
  \textbf{Repository} & \href{https://github.com/Graf-J/MLDS-Q4-490-Advanced-Algorithms-HW2}{\texttt{github.com/Graf-J/MLDS-Q4-490-Advanced-Algorithms-HW2}} (public)
\end{tabularx}

The project uses \href{https://docs.astral.sh/uv/}{uv} for Python and its dependencies (PyTorch with CUDA):

\begin{lstlisting}
git clone https://github.com/Graf-J/MLDS-Q4-490-Advanced-Algorithms-HW2.git
cd MLDS-Q4-490-Advanced-Algorithms-HW2
uv sync             # creates .venv with the exact versions from uv.lock
uv run jupyter lab  # or open the notebooks in VS Code with the .venv kernel
\end{lstlisting}

Every part of the assignment is one Jupyter notebook that runs from top to bottom. All figures are written to \texttt{report/figures/}.

\begin{center}
\small
\begin{tabularx}{\textwidth}{@{}l l L@{}}
  \toprule
  \textbf{Part} & \textbf{Notebook} & \textbf{Content} \\
  \midrule
  a: CartPole & \texttt{cartpole-v1.ipynb} & Hyperparameter ladder, seed study, phase portraits, edge cases and validation tests (Section~\ref{sec:cartpole}) \\
  b: Ms.\ Pac-Man & \texttt{mspacman-v5.ipynb} & Training with crash recovery, model selection, evaluation, heatmaps, edge cases and validation tests (Section~\ref{sec:mspacman}); the last cell opens a window in which the agent plays \\
  \bottomrule
\end{tabularx}
\end{center}

\begin{sloppypar}
\textbf{Running without retraining.} Every training run is saved to \texttt{checkpoints/a/<variant>/seed<k>/} (training log, network weights, config). With the default flag \texttt{RETRAIN = False}, the notebook loads a saved run instead of training it again, as long as its saved config equals the current one. A changed config is retrained automatically. All results in this report can therefore be reproduced from the repository without any training.

For part (b), the repository contains the training log, the network snapshots (every 500 episodes) and the cached evaluations in \texttt{checkpoints/b/dqn/}, but not the full checkpoint with the 7\,GB replay buffer. The notebook recognises the finished run and loads it instead of training again. A new training (about 6 hours on a laptop GPU, 32\,GB of RAM recommended) resumes automatically from the last checkpoint if it is interrupted.
\end{sloppypar}

\begin{notebox}
\textbf{How to run the validation tests.} Every test described in this report is implemented in Section~7 \emph{Validation Tests} of its notebook (part (a): Tests 1--9, part (b): Tests 1--10), each as a markdown cell followed by its code:
\begin{itemize}
  \item The markdown cell explains \textbf{what} the test does and its \textbf{pass criterion}; why each test validates the algorithm is discussed in this report.
  \item The code cell runs the experiment and prints \texttt{[PASS]} or \texttt{[FAIL]} together with the measured numbers. A summary cell at the end of the notebook lists all checks.
  \item The tests run automatically with \emph{Run All}. To rerun a single test, run all cells before Section~7 (they load the saved runs instead of training) and then the test's cells.
  \item Tests that need extra trainings: the control experiments and ablations of part (a) are saved in \texttt{checkpoints/a/} and loaded like all other runs. The $\gamma = 0$ run and the resume test of part (b) train small runs of a few minutes. Expensive evaluations of part (b) are cached in \texttt{checkpoints/b/dqn/cache/}; deleting this folder recomputes them.
\end{itemize}
In this report, a passed check is marked with~\pass{} and a failed one with~\fail.
\end{notebox}

\textbf{Hardware.} The networks were trained on a single GPU (NVIDIA GeForce GTX 1650 Max-Q, 4\,GB) in a laptop with 32\,GB of RAM.
